\chapter{Dalsze narzędzia homotopii: granice odwrotne, przeszkody, operacje i sklejanie Morse'a}
\label{ch:xviii-domkniecie}

\section{Cel, motywacja i poziom dowodów}
W poprzednich rozdziałach spotkaliśmy cztery zjawiska wymagające
szczególnej ostrożności: przechodzenie od skończonych przybliżeń do
przestrzeni granicznej; rozszerzanie przekrojów krok po kroku; zależności
między operacjami kohomologicznymi; oraz sklejanie trajektorii przy
złamanych liniach gradientowych. Łączy je pytanie, czy dane zgodne na
każdym skończonym etapie dają obiekt globalny. Odpowiedź może być
negatywna nawet wtedy, gdy wszystkie lokalne grupy przeszkód wydają się
proste.

Celem rozdziału jest ścisłe uzupełnienie tych czterech zagadnień.
Dowodzimy elementarnej teorii $\varprojlim^1$ dla przeliczalnych
układów grup abelowych i odpowiedniego ciągu Milnora dla układów
kompleksów łańcuchowych. Przeprowadzamy konstrukcję kolejnych
przeszkód komórkowych, wraz z kontrolą wpływu wyborów. Dla operacji
Steenroda wyprowadzamy kilka relacji z Bocksteina i podajemy dokładny
zakres twierdzenia Adema. Dla ekscyzji homotopijnej wyjaśniamy
mechanizm twierdzenia Blakersa--Masseya i jego zastosowanie do
zawieszenia. W końcu identyfikujemy precyzyjnie dodatkowy lemat
analityczny potrzebny do pełnego sklejania trajektorii Morse'a.

\section{Układy odwrotne i dwa rodzaje granic}
\begin{definicja}[Przeliczalny układ odwrotny]
Przez układ odwrotny grup abelowych rozumiemy ciąg
$A_0\xleftarrow{u_0}A_1\xleftarrow{u_1}A_2\leftarrow\cdots$,
gdzie każdy $A_j$ jest grupą abelową, a $u_j:A_{j+1}\to A_j$
homomorfizmem. Jego granicę odwrotną definiujemy jako
\[
\varprojlim_jA_j=
\{(a_j)_{j\ge0}\in\prod_{j\ge0}A_j:
 a_j=u_j(a_{j+1})\text{ dla każdego }j\}.
\]
Sytuację przedstawia schemat odwzorowań przejścia:
\begin{center}
\begin{tikzpicture}[>=Stealth, node distance=2.0cm,
  every node/.style={font=\small}]
\node (a0) {$A_0$};
\node[right=of a0] (a1) {$A_1$};
\node[right=of a1] (a2) {$A_2$};
\node[right=of a2] (a3) {$\cdots$};
\draw[->] (a1) -- node[above] {$u_0$} (a0);
\draw[->] (a2) -- node[above] {$u_1$} (a1);
\draw[->] (a3) -- node[above] {$u_2$} (a2);
\end{tikzpicture}
\par\smallskip\small Rysunek: strzałki układu odwrotnego wskazują od dokładniejszych przybliżeń do wcześniejszych.
\end{center}
Na produkcie definiujemy operator
$T((a_j)_{j\ge0})=(u_j(a_{j+1}))_{j\ge0}$.
\end{definicja}
\begin{definicja}[Pierwszy pochodny funktor granicy odwrotnej]
Niech $(A_j,u_j)_{j\ge0}$ będzie przeliczalnym
układem odwrotnym grup abelowych i niech
$T:\prod_{j\ge0}A_j\to\prod_{j\ge0}A_j$
będzie operatorem
$T((a_j))=(u_j(a_{j+1}))_{j\ge0}$.
Kładziemy
\[
\varprojlim{}^1_jA_j:=
\operatorname{coker}\bigl(1-T:\prod_{j\ge0}A_j\to
\prod_{j\ge0}A_j\bigr).
\]
Nie jest to nowy typ granicy przestrzeni topologicznych:
$\varprojlim{}^1$ jest grupą mierzącą brak dokładności
funktora $\varprojlim$ na takich układach.
W kategorii przeliczalnych układów grup abelowych
iloczyny grup są dokładne, a powyższy kokernel
realizuje pierwszy prawy funktor pochodny granicy odwrotnej.
\end{definicja}
\begin{lemat}[Jądro operatora różnicy]
Niech $(A_j,u_j)_{j\ge0}$ będzie przeliczalnym
układem odwrotnym grup abelowych i
$T((a_j))=(u_j(a_{j+1}))$ na $\prod_{j\ge0}A_j$.
Wówczas zachodzi
$\ker(1-T)=\varprojlim_j A_j$.
\end{lemat}
\begin{proof}
Równość $(1-T)(a)=0$ oznacza dokładnie układ równań
$a_j=u_j(a_{j+1})$ dla wszystkich $j$. Są to wszystkie i tylko
warunki definiujące granicę odwrotną.
\end{proof}

\begin{twierdzenie}[Warunek Mittaga--Lefflera dla układów przeliczalnych]
\label{thm:xviii-ml}
Niech $(A_j,u_j)$ będzie przeliczalnym układem grup abelowych.
Załóżmy, że dla każdego $j$ obrazy
$\operatorname{im}(A_k\to A_j)$ stabilizują się, gdy $k\to\infty$.
Wówczas $\varprojlim{}^1 A_j=0$.
W szczególności wniosek zachodzi, gdy wszystkie $u_j$ są surjektywne.
\end{twierdzenie}
\begin{proof}[Przypadek surjektywnych odwzorowań]
Pokażemy, że $1-T$ jest surjektywny. Ustalmy
$(b_j)\in\prod_jA_j$. Chcemy rozwiązać
$a_j-u_j(a_{j+1})=b_j$ dla wszystkich $j$.
Wybierzmy $a_0=0$. Jeśli $a_j$ jest już wybrany, to
surjektywność $u_j$ pozwala wybrać $a_{j+1}$ z równania
$u_j(a_{j+1})=a_j-b_j$. Indukcja daje szukany ciąg,
zatem kokernel $1-T$ jest zerowy.
\end{proof}
\begin{proof}[Redukcja ogólnego przypadku]
Dla każdego $j$ niech $B_j\subset A_j$ będzie ustabilizowanym
obrazem $A_k\to A_j$ dla dostatecznie dużych $k$.
Odwzorowania $B_{j+1}\to B_j$ są surjektywne: jeśli
$b\in B_j$, wybieramy $k$ większe od progów stabilizacji dla
$j$ i $j+1$ i element $a_k$ przechodzący na $b$. Jego obraz
w $B_{j+1}$ przechodzi na $b$. Zatem $\varprojlim{}^1 B_j=0$.

Dla układu ilorazów $Q_j=A_j/B_j$ każdy element $Q_k$
przechodzi na zero w $Q_j$ dla dostatecznie dużego $k$.
Taki układ nazywamy pro-zerowym. Dla ustalonego ciągu
$(q_j)\in\prod_j Q_j$ rozważmy równania
$x_j-v_j(x_{j+1})=q_j$. Możemy je rozwiązać przez sumę
skończoną w każdej współrzędnej:
$x_j=\sum_{k\ge j}v_{j,k}(q_k)$, gdzie
$v_{j,k}:Q_k\to Q_j$ jest odwzorowaniem przejścia
(i $v_{j,j}=\id$). Suma jest rzeczywiście skończona,
bo dla ustalonego $j$ odwzorowania $v_{j,k}$ zerują się dla dużych $k$.
Wtedy $x_j-v_j(x_{j+1})=q_j$. Stąd $\varprojlim{}^1Q_j=0$.

Pozostaje użyć dokładności sześciowyrazowego ciągu dla
krótkiego ciągu układów $0\to B_j\to A_j\to Q_j\to0$.
Dokładniej, lemat o wężu dla dwóch dokładnych
wierszy produktów i pionowych odwzorowań $1-T$
daje fragment ciągu dokładnego
\[
\varprojlim{}^1 B_j\longrightarrow
\varprojlim{}^1 A_j\longrightarrow
\varprojlim{}^1 Q_j\longrightarrow0.
\]
Pierwszy wyraz znika z powodu surjektywności
$B_{j+1}\to B_j$, a ostatni z powodu pro-zerowości
układu $Q_j$. Stąd $\varprojlim{}^1A_j=0$.
\end{proof}
\begin{uwaga}
W powyższym dowodzie użyto produktów grup abelowych,
które zachowują dokładność krótkich ciągów. W ogólniejszych
kategoriach homologicznym warunkiem zerowania mogą być inne
hipotezy; symbol $\varprojlim{}^1$ nie jest wtedy automatycznie
definiowany przez ten sam elementarny iloraz.
\end{uwaga}

\section{Przykład podstawowy: $\mathbb Z_p/\mathbb Z$}
\begin{definicja}
Niech $p\ge2$ będzie liczbą pierwszą. Pierścieniem liczb
$p$-adycznych nazywamy granicę odwrotną
$\mathbb Z_p=\varprojlim_m\mathbb Z/p^m\mathbb Z$,
gdzie odwzorowaniami przejścia są redukcje modulo $p^m$.
\end{definicja}
\begin{twierdzenie}[Niezerowe $\varprojlim{}^1$]
\label{thm:xviii-padic}
Niech $p$ będzie liczbą pierwszą,
$A_j=\Z$ dla wszystkich $j\ge0$ oraz
$u_j:A_{j+1}\to A_j$ homomorfizmem
$u_j(a)=pa$. Wówczas zachodzą
izomorfizmy
\[
\varprojlim_j A_j=0,\qquad
\varprojlim{}^1_j A_j\cong\Z_p/\Z,
\]
gdzie $\Z$ zanurzamy w $\Z_p$ przez zwykłe redukcje
modulo potęg $p$.
\end{twierdzenie}
\begin{proof}
Jeśli $a_j=pa_{j+1}$ dla wszystkich $j$, to $a_0$
jest podzielne przez $p^m$ dla każdego $m$, zatem $a_0=0$.
Indukcyjnie otrzymujemy $a_j=0$.

Niech $P=\prod_{j\ge0}\Z$. Definiujemy homomorfizm
$\Theta:P\to\Z_p$ przez zbieżny szereg $p$-adyczny
$\Theta((a_j))=\sum_{j\ge0}p^ja_j$.
Surjektywność wynika z rozwinięcia każdej liczby $p$-adycznej
w cyfry $0,\ldots,p-1$.
Jeżeli $b=(1-T)a$, czyli $b_j=a_j-pa_{j+1}$,
to sumy częściowe teleskopują:
$\sum_{j=0}^{m}p^jb_j=a_0-p^{m+1}a_{m+1}$.
Ostatni wyraz dąży do zera $p$-adycznie, więc
$\Theta(b)=a_0\in\Z$. Dostajemy indukowane odwzorowanie
$\overline\Theta:\operatorname{coker}(1-T)\to\Z_p/\Z$.

Aby wykazać różnowartościowość, załóżmy, że
$\Theta(b)=q\in\Z$. Ustalmy $s_0=q$ i dla $j\ge0$ kładźmy
\[
s_j=\frac{q-\sum_{i=0}^{j-1}p^ib_i}{p^j}.
\]
Licznik jest podzielny przez $p^j$ w $\Z_p$, a ponieważ
jest liczbą całkowitą, jest podzielny przez $p^j$ w $\Z$.
Zatem $s_j\in\Z$ i $s_j-ps_{j+1}=b_j$.
Stąd $b=(1-T)s$, co dowodzi różnowartościowości.
Surjektywność $\overline\Theta$ jest dziedziczona po $\Theta$.
\end{proof}
\begin{przyklad}[Skończona rozwiązywalność bez rozwiązania globalnego]
\label{prz:xviii-lim1-2adic-audit}
W Twierdzeniu~\ref{thm:xviii-padic} weźmy $p=2$
i ciąg $b=(b_j)_{j\ge0}$ określony przez
$b_{2m}=1$ oraz $b_{2m+1}=0$ dla $m\ge0$.
Rozpatrujemy układ nieskończenie wielu równań
\[
 a_j-2a_{j+1}=b_j,\qquad j=0,1,2,\ldots,
 \qquad a_j\in\Z.
\]
\emph{Każdy skończony fragment ma rozwiązanie.}
Dla dowolnego $N\ge1$ połóżmy $a_N=0$ i wyznaczmy
wstecz $a_j=b_j+2a_{j+1}$ dla $j=N-1,\ldots,0$.
Otrzymujemy liczby całkowite spełniające
wszystkie równania o indeksach $0,\ldots,N-1$.

\emph{Nie ma rozwiązania wszystkich równań naraz.}
Gdyby $(a_j)_{j\ge0}$ było takim rozwiązaniem,
to z rachunku teleskopowego w dowodzie
Twierdzenia~\ref{thm:xviii-padic} mielibyśmy
\[
 a_0=\Theta(b)=\sum_{m=0}^{\infty}2^{2m}
 =\sum_{m=0}^{\infty}4^m\qquad\text{w }\Z_2.
\]
Szereg jest zbieżny $2$-adycznie, ponieważ
$4^m\to0$ w tej topologii. Aby skontrolować
każdy etap, dla $M\ge0$ obliczmy w $\Z$
$3\sum_{m=0}^M4^m=4^{M+1}-1$.
Stąd sumy częściowe są zgodne z $-1/3$
modulo $2^{2M+2}$, co uzasadnia również
kolejne redukcje przedstawione na rysunku poniżej.
Dla sumy $S$ otrzymujemy $4S=S-1$,
czyli $3S=-1$.
Zatem $S=-1/3\in\Z_2$, gdyż $3$ jest
odwracalne w pierścieniu liczb $2$-adycznych.
Żadna liczba całkowita $a_0$ nie spełnia
$3a_0=-1$; sprzeczność.
Klasa $b$ jest więc \emph{niezerowa} w
$\varprojlim{}^1(\Z,\times2)\cong\Z_2/\Z$.
To konkretny przykład przeszkody globalnej,
mimo że na każdym skończonym etapie można
dobrać rozwiązania.
\end{przyklad}

\begin{figure}[htbp]
\centering
\begin{tikzpicture}[>=Stealth,node distance=12mm and 13mm,
  box/.style={draw,rounded corners=3pt,inner sep=5pt,font=\small}]
\node[box] (a) {$r_2=1\pmod4$};
\node[box,right=of a] (b) {$r_3=5\pmod8$};
\node[box,right=of b] (c) {$r_4=5\pmod{16}$};
\node[box,right=of c] (d) {$r_5=21\pmod{32}$};
\draw[->,thick,manifoldblue] (b)--(a);
\draw[->,thick,manifoldblue] (c)--(b);
\draw[->,thick,manifoldblue] (d)--(c);
\end{tikzpicture}
\caption{Zgodne redukcje początkowych sum szeregu
$1+4+16+\cdots$ w $\Z_2$.
Granica jest liczbą $-1/3\in\Z_2\setminus\Z$.
Strzałki oznaczają redukcję modulo mniejszą potęgę $2$.}
\label{fig:xviii-lim1-2adic-audit}
\end{figure}

\begin{przyklad}[Dlaczego granicy nie można przestawiać z homologią]
Przyjmijmy $p=2$. Dla $j\ge0$ weźmy
kompleks $C^{(j)}_*$ skupiony w stopniach $1$ i $0$,
o postaci
\[
C^{(j)}_1=\Z\xrightarrow{\ \times p^j\ }
C^{(j)}_0=\Z.
\]
Określmy odwzorowanie łańcuchowe
$u_j:C^{(j+1)}_*\to C^{(j)}_*$
przez mnożenie przez $p$ w stopniu $1$ i identyczność
w stopniu $0$. Rzeczywiście,
$(\times p^j)\circ(\times p)
=\id\circ(\times p^{j+1})$,
więc otrzymujemy układ odwrotny kompleksów.
Jego granica w stopniu $1$ jest zerowa:
ciąg liczb całkowitych spełniający
$a_j=pa_{j+1}$ musi być zerowy.
W stopniu $0$ odwzorowania przejścia są
identycznościami, a zatem granica wynosi $\Z$.
Dostajemy $H_0(\varprojlim_j C^{(j)}_*)\cong\Z$.
Z drugiej strony
$H_0(C^{(j)}_*)\cong\Z/p^j\Z$,
a indukowane odwzorowania przejścia są
redukcjami, więc
$\varprojlim_j H_0(C^{(j)}_*)\cong\Z_p$.
Naturalne odwzorowanie $\Z\to\Z_2$ nie jest
surjektywne: z Przykładu~\ref{prz:xviii-lim1-2adic-audit}
wiemy, że $-1/3\in\Z_2\setminus\Z$.
W szczególności homologia i granica odwrotna
nie są tutaj przemienne.

Nie jest to kontrprzykład do
Twierdzenia~\ref{thm:xviii-milnor}:
odwzorowania $u_j$ w stopniu $1$
są mnożeniem przez $p$ na $\Z$ i \emph{nie}
są surjektywne. Surjektywność odwzorowań
na homologii nie zastępuje surjektywności
w każdym stopniu łańcuchowym.
\end{przyklad}

\section{Ciąg Milnora dla układów kompleksów}
\begin{twierdzenie}[Elementarny ciąg Milnora]
\label{thm:xviii-milnor}
Niech $\{C^{(j)}_*,u_j\}_{j\ge0}$ będzie przeliczalnym
układem odwrotnym kompleksów łańcuchowych grup abelowych,
a każde $u_j:C^{(j+1)}_q\to C^{(j)}_q$ niech będzie
surjektywne dla wszystkich $j,q$. Wówczas dla każdego $n\in\Z$
istnieje naturalny krótki ciąg dokładny
\[
0\longrightarrow\varprojlim{}^1_jH_{n+1}(C^{(j)}_*)
\longrightarrow H_n(\varprojlim_j C^{(j)}_*)
\longrightarrow\varprojlim_jH_n(C^{(j)}_*)
\longrightarrow0.
\]
Na ogół ten ciąg nie ma naturalnego rozszczepienia.
\end{twierdzenie}
\begin{proof}
Oznaczmy przez $P_*$ produkt kompleksów $\prod_j C^{(j)}_*$,
z różniczką działającą współrzędnie. Odwzorowania łańcuchowe $u_j$
definiują endomorfizm $T:P_*\to P_*$. Surjektywność map
$u_j$ i argument indukcyjny użyty w twierdzeniu
\ref{thm:xviii-ml} pokazują, że $1-T:P_q\to P_q$
jest surjektywny w każdym stopniu. Jego jądro to z definicji
$\varprojlim C^{(j)}_q$. Otrzymujemy krótki ciąg dokładny
kompleksów $0\to\varprojlim C^{(j)}_*\to P_*
\xrightarrow{1-T}P_*\to0$.

Długi ciąg homologii zawiera fragment
\[
H_{n+1}(P_*)\xrightarrow{1-T}H_{n+1}(P_*)
\longrightarrow H_n(\varprojlim C^{(j)}_*)
\longrightarrow H_n(P_*)\xrightarrow{1-T}H_n(P_*).
\]
Ponieważ w kategorii grup abelowych produkty zachowują
jądra i obrazy różniczek, mamy
$H_q(P_*)\cong\prod_jH_q(C^{(j)}_*)$.
Jądro prawego odwzorowania to $\varprojlim H_n(C^{(j)}_*)$,
a kokernel lewej to $\varprojlim{}^1 H_{n+1}(C^{(j)}_*)$.
Wyodrębnienie tego fragmentu daje żądany krótki ciąg.
\end{proof}
\begin{uwaga}[Czego twierdzenie nie mówi]
Nie wolno zastosować ciągu~\ref{thm:xviii-milnor}
automatycznie do $C_*(\varprojlim X_j)$, nawet gdy
$X_j$ są zwartymi wielościanami. Zwykle brak naturalnego
izomorfizmu $C_*(\varprojlim X_j)\cong\varprojlim C_*(X_j)$,
a odwzorowania na grupach łańcuchów nie muszą być surjektywne.
Jest to jedna z przyczyn, dla których homologia Čecha i homologia
singularna zachowują się odmiennie na przestrzeniach dzikich.
\end{uwaga}

\section{Wyższe przeszkody dla przekrojów i odwzorowań}
\begin{definicja}[Problem rozszerzania przekroju]
Niech $p:E\to B$ będzie fibracją Serre'a nad spójnym kompleksem CW,
niech $B^{(q)}$ oznacza jego $q$-szkielet, a $s_q:B^{(q)}\to E$
ciągły przekrój nad $B^{(q)}$. Problem polega na rozszerzeniu
$s_q$ na $B^{(q+1)}$, a ostatecznie na całe $B$.
Zakładamy, że włókno $F$ jest łukowo spójne; przy formułowaniu
przeszkód z abelowymi współczynnikami zakładamy też $q\ge2$.
\end{definicja}
\begin{stwierdzenie}[Kołańcuch przeszkody]
\label{prop:xviii-obstruction}
Niech $q\ge2$, niech $p:E\to B$ będzie fibracją
Serre'a nad spójnym kompleksem CW $B$ z ustalonym
punktem bazowym, a jej włókno $F$ niech będzie
$(q-1)$-spójne. Niech
$s_q:B^{(q)}\to E$ będzie ciągłym przekrojem
nad $q$-szkieletem. Niech
$\underline{\pi_q(F)}$ oznacza lokalny system
grup abelowych wyznaczony przez transport
homotopijny włókien fibracji.
Dla każdej zorientowanej komórki $e^{q+1}$ wybierzmy
jej charakterystykę $\chi_e:D^{q+1}\to B$ i trywializację
ściągniętej fibracji nad dyskiem w sensie równoważności
włóknowej z fibracją produktową, wraz z wyborem transportów
identyfikujących włókna. Nie zakładamy, że trywializacja
jest już zadana i ustalona na całym brzegu dysku:
to właśnie możliwość przedłużenia przekroju przez dysk
jest przedmiotem badanej przeszkody.
Klasa ograniczenia
$s_q\circ\chi_e|_{S^q}$ określa element $\pi_q(F)$.
Po uwzględnieniu transportu wzdłuż $1$-szkieletu otrzymujemy
kołańcuch komórkowy
$o_{q+1}(s_q)\in C^{q+1}(B;\underline{\pi_q(F)})$.
Jest on kocyklem, a jego klasa w
$H^{q+1}(B;\underline{\pi_q(F)})$ jest pierwszą przeszkodą
rozszerzenia przekroju. Jeśli ta klasa znika, można zmienić
przekrój na $q$-szkielecie i rozszerzyć go na $B^{(q+1)}$.
\end{stwierdzenie}
\begin{proof}[Konstrukcja i kontrola zależności]
Na każdej komórce dyskowej pullback fibracji jest
włóknowo homotopijnie trywialny: stosujemy własność
podnoszenia homotopii do ściągnięcia dysku do punktu.
Dzięki $(q-1)$-spójności $F$ wszystkie rozszerzenia do
$q$-szkieletu można porównywać, a element otrzymany na
brzegu komórki jest dobrze określony po wybraniu transportów
między włóknami. Monodromia $\pi_1(B)$ może działać
nietrywialnie na $\pi_q(F)$; właśnie to działanie definiuje
system lokalnych współczynników podkreślony w tezie.

Warunek kocyklu można sprawdzić na jednej komórce
$e^{q+2}$. Jej przyczepienie daje sumę granicznych map na
$q$-sferach ścian, przemnożonych przez odpowiednie
homomorfizmy transportu. Suma tych klas jest zerowa,
ponieważ jest obrazem brzegu odwzorowania dysku po przejściu do
łańcuchów komórkowych i wykorzystaniu zgodności
przekształceń homotopii na przecięciach. Ten argument
jest standardowym \emph{lematem komórkowym o kocyklu
przeszkody}; jego pełna konstrukcja wymaga ułożenia
homotopii charakterystyk na $(q+2)$-komórkach.

Jeżeli zmienimy $s_q$ na $q$-komórkach przez klasy
$\beta\in C^q(B;\underline{\pi_q(F)})$, kołańcuch
przeszkody zmienia się o $\delta\beta$ (z odpowiednią
konwencją znaków). Zatem klasa kohomologii jest
niezależna od tej zmiany. Jej znikanie umożliwia wybór
$\beta$ zabijający kołańcuch, po czym każde odwzorowanie
$S^q\to F$ z trywialną klasą homotopii rozszerza się
na $D^{q+1}$. Z własności podnoszenia otrzymujemy
przekrój nad każdym dołączonym dyskiem i ich zgodność
na brzegu. To kończy argument pod warunkiem przyjęcia
wskazanego lematu komórkowego, którego nie zastępujemy
samym sloganem $\partial^2=0$.
\end{proof}
\begin{uwaga}[Kolejne przeszkody]
Po zniwelowaniu pierwszej klasy można przejść do następnego
szkieletu. Następna przeszkoda leży w grupie z wyższymi
homotopiami włókna, ale \emph{zależy od dokonanego wyboru
wcześniejszego rozszerzenia}. Nie ma na ogół jednej
kanonicznej listy klas niezależnych od siebie.
Ponadto przy $q=1$ współczynniki mogą być nieabelowe;
wtedy klasyczny zapis grupy kohomologii wymaga odrębnych
konstrukcji. W przypadku $q\ge2$ skręcenie współczynników
jest wyznaczone przez akcję monodromii na $\pi_q(F)$.
\end{uwaga}
\begin{przyklad}[Wiązka sferowa i klasa Eulera]
Dla zorientowanej rzeczywistej wiązki wektorowej $\xi$ rzędu
$r\ge3$ nad spójnym CW-kompleksem $B$ włóknem wiązki
jednostkowych sfer jest $S^{r-1}$, które jest $(r-2)$-spójne.
Pierwsza przeszkoda istnienia przekroju nigdzie niezerowego
leży w $H^r(B;\Z)$ i jest klasą Eulera $e(\xi)$.
Jeżeli $\dim B=r$, jej znikanie jest także wystarczające,
gdyż nie ma komórek wyższych stopni. Jeżeli
$\dim B>r$, mogą występować wyższe przeszkody, czego
nie wykrywa sama klasa Eulera.
\end{przyklad}
\begin{przyklad}[Zerowa klasa Eulera nie wystarcza]
Dla zorientowanej wiązki rzeczywistej rzędu trzy nad
$S^4$ klasa Eulera należy do $H^3(S^4;\Z)=0$.
Mimo to istnieją takie wiązki bez przekroju nigdzie
niezerowego: nietrywialne elementy
$\pi_3(SO(3))\cong\Z$ wyznaczają przez sklejanie
wiązki rzędu trzy. Przekrój wymuszałby redukcję
struktury do $SO(2)$, a wtedy klasa sklejenia pochodziłaby
z $\pi_3(SO(2))=0$.
Przykład pokazuje rzeczywistą wyższą przeszkodę.
\end{przyklad}

\section{Iloczyny Steenroda i relacje Adema: zakres elementarny}
Na kohomologii z $\F_2$ działają naturalne operacje
$\operatorname{Sq}^i:H^n(X;\F_2)\to H^{n+i}(X;\F_2)$.
Przypominamy normalizację:
$\operatorname{Sq}^0=\id$,
$\operatorname{Sq}^n(x)=x\smile x$ dla $x\in H^n$,
$\operatorname{Sq}^i(x)=0$ dla $i>n$.
Do tego dochodzą naturalność i wzór Cartana, już omówione
we wcześniejszych rozdziałach.
\begin{twierdzenie}[Relacje Adema, $p=2$]
\label{thm:xviii-adem}
Niech $a,b$ będą liczbami całkowitymi,
$0<a<2b$. W kohomologii ze współczynnikami
w ciele $\F_2$, dla wszystkich przestrzeni,
na których zdefiniowane są kwadraty Steenroda,
obowiązuje równość operacji naturalnych
\[
\operatorname{Sq}^a\operatorname{Sq}^b=
\sum_{j=0}^{\lfloor a/2\rfloor}
\binom{b-j-1}{a-2j}
\operatorname{Sq}^{a+b-j}\operatorname{Sq}^{j},
\]
przy czym współczynniki dwumianowe liczymy modulo dwa.
Wzór jest twierdzeniem o operacjach na \emph{wszystkich}
przestrzeniach, nie tylko o ich wartościach w jednym pierścieniu
kohomologii.
\end{twierdzenie}
\begin{uwaga}[Status dowodu]
Pełny dowód wymaga kołańcuchowej homotopii relacji Adema
lub uniwersalnej konstrukcji operacji i obliczeń w
$H^*(K(\F_2,n);\F_2)$. Nie wynika on wyłącznie z
naturalności i wzoru Cartana. W tym rozdziale przedstawiamy
jedynie samodzielne wyprowadzenia kilku konsekwencji i
kontrolę obliczeniową; twierdzenie~\ref{thm:xviii-adem}
przyjmujemy jako zewnętrzny wynik.
\end{uwaga}
\begin{stwierdzenie}[Pierwsze relacje]
Z relacji Adema wynikają
$\operatorname{Sq}^1\operatorname{Sq}^1=0$,
$\operatorname{Sq}^1\operatorname{Sq}^2=\operatorname{Sq}^3$,
$\operatorname{Sq}^2\operatorname{Sq}^2=
\operatorname{Sq}^3\operatorname{Sq}^1$.
\end{stwierdzenie}
\begin{proof}
Podstawiamy w\,~\ref{thm:xviii-adem} odpowiednio
$(a,b)=(1,1),(1,2),(2,2)$. W ostatnim przypadku
składnik $j=0$ ma współczynnik
$\binom{1}{2}=0$, a składnik $j=1$ ma
$\binom{0}{0}=1$. Dla $(1,1)$ jedyny współczynnik to
$\binom{0}{1}=0$; dla $(1,2)$ to $\binom{1}{1}=1$.
\end{proof}
\begin{przyklad}[Obliczenie na $\mathbb RP^{\infty}$]
Dla generatora $t\in H^1(\mathbb RP^\infty;\F_2)$ mamy
$\operatorname{Sq}^i(t^m)=\binom mi t^{m+i}$.
Wzór wynika z $\operatorname{Sq}(t)=t+t^2$ i
wzoru Cartana: $\operatorname{Sq}(t^m)=(t+t^2)^m$.
Na przykład
$\operatorname{Sq}^2\operatorname{Sq}^2(t^3)=0$,
a $\operatorname{Sq}^3\operatorname{Sq}^1(t^3)=0$.
Jest to zgodność z relacją, ale samo obliczenie na tej
przestrzeni \emph{nie dowodzi} relacji dla wszystkich $X$.
\end{przyklad}
\begin{przyklad}[Przypadek $\mathbb CP^\infty$]
Jeśli $u\in H^2(\mathbb CP^\infty;\F_2)$ jest redukcją
pierwszej klasy Cherna uniwersalnej wiązki liniowej, to
$\operatorname{Sq}^2(u)=u^2$ i
$\operatorname{Sq}^{2i}(u^m)=\binom mi u^{m+i}$,
podczas gdy nieparzyste kwadraty znikają na całym
$\F_2[u]$ z przyczyn parzystości stopni.
\end{przyklad}

\section{Ekscyzja homotopijna Blakersa--Masseya}
\begin{definicja}[Spójność pary]
Para przestrzeni z punktem bazowym $(X,A)$ jest $m$-spójna,
$m\ge1$, jeżeli $A$ i $X$ są łukowo spójne i
$\pi_j(X,A)=0$ dla wszystkich $1\le j\le m$.
Przyjmiemy typowe hipotezy CW, aby grupy względne i
podnoszenie homotopii miały oczekiwane własności.
\end{definicja}
\begin{twierdzenie}[Blakers--Massey, postać ekscyzyjna]
\label{thm:xviii-bm}
Niech $X=A\cup B$ będzie triadą CW, a $C=A\cap B$ niech
będzie niepustym łukowo spójnym podkompleksem.
Załóżmy, że $(A,C)$ jest $m$-spójna, $(B,C)$ jest
$n$-spójna, gdzie $m,n\ge1$.
Wówczas homomorfizm ekscyzji
\[
\pi_q(A,C)\longrightarrow\pi_q(X,B)
\]
jest izomorfizmem dla $q<m+n$ i epimorfizmem dla $q=m+n$.
W stopniu $q=1$ rozumiemy odpowiednie zbiory punktowane
lub grupy względne zgodnie ze zwykłą konwencją.
\end{twierdzenie}
\begin{uwaga}[Zakres dowodu i źródło trudności]
Dla homologii ekscyzja jest izomorfizmem bez ograniczenia
stopnia. W homotopii zakres zależy od spójności obu par.
Dowód Blakersa--Masseya polega na metodzie \emph{usuwania
przecięć}: reprezentant odwzorowania dysku do $X$ rozcina
się na obszary przechodzące do $A$ i $B$, a założenia
spójności pozwalają sukcesywnie zmieniać reprezentanta,
usuwając obszary niewłaściwe w zakresie wymiarów
$q\le m+n$. Ten geometryczny argument wymaga
technicznego lematu o uporządkowaniu przecinających się
komórek; nie uzasadnia go sama ekscyzja homologiczna.
Pełny lemat geometryczny przyjmujemy jako twierdzenie
zewnętrzne. Poniżej dowodzimy wynikający z niego
zakres stabilizacji Freudenthala.
\end{uwaga}
\begin{wniosek}[Stabilizacja Freudenthala]
\label{cor:xviii-freudenthal}
Niech $Y$ będzie $(r-1)$-spójnym, łukowo spójnym
kompleksem CW, gdzie $r\ge2$. Wówczas homomorfizm
zawieszenia
$\pi_k(Y)\to\pi_{k+1}(\Sigma Y)$ jest izomorfizmem
dla $k\le2r-2$ i epimorfizmem dla $k=2r-1$.
\end{wniosek}
\begin{proof}
Rozłóżmy $\Sigma Y$ na dwa stożki $A=CY$ i $B=CY$,
których przecięcie jest homotopijnie równoważne $Y$.
Po odpowiednim zastąpieniu przez triadę CW
obie pary $(CY,Y)$ są $r$-spójne, ponieważ
$\pi_j(CY,Y)\cong\pi_{j-1}(Y)$ dla $j\ge2$.
Twierdzenie~\ref{thm:xviii-bm} daje izomorfizm
$\pi_q(CY,Y)\to\pi_q(\Sigma Y,CY)$ dla $q<2r$
i epimorfizm dla $q=2r$.
Długie ciągi dokładne par oraz ściągalność stożków
dają $\pi_q(CY,Y)\cong\pi_{q-1}(Y)$ i
$\pi_q(\Sigma Y,CY)\cong\pi_q(\Sigma Y)$ dla $q\ge2$.
Przy tych identyfikacjach odwzorowanie ekscyzji odpowiada
homomorfizmowi zawieszenia, co daje wynik po podstawieniu
$q=k+1$.
\end{proof}
\begin{przyklad}[Zawieszenie okręgu]
Dla $Y=S^1$ powyższy wniosek przy $r=1$ nie ma zastosowania
w tej postaci; można jednak bezpośrednio obliczyć
$\Sigma S^1\cong S^2$ oraz
$\pi_1(S^1)\cong\Z$, $\pi_2(S^2)\cong\Z$.
Ważne jest niewnioskowanie ogólnego zakresu stabilizacji
z jednego przykładu.
\end{przyklad}

\section{Sklejanie trajektorii Morse'a: ścisły bilans analityczny}
Niech $M$ będzie zwartą gładką rozmaitością bez brzegu,
$f:M\to\R$ funkcją Morse'a, a $g$ metryką Riemannowską
taką, że para $(f,g)$ spełnia warunek Morse'a--Smale'a.
Dla punktów krytycznych $p,q$ oznaczmy przez
$\mathcal M(p,q)$ przestrzeń trajektorii równania
$\dot\gamma=-\nabla_g f(\gamma)$, modulo przesunięcia czasu.
Przyjmujemy $\operatorname{ind}(p)-\operatorname{ind}(q)=2$.
Chcemy opisać końce jednowymiarowej przestrzeni
$\mathcal M(p,q)$ jako złamane trajektorie przechodzące
przez punkty o indeksie pośrednim.

\begin{definicja}[Operator linearyzacji i norma ważona]
Niech $\gamma:\R\to M$ będzie trajektorią gradientową.
Wzdłuż trajektorii $\gamma$ używamy pochodnej
kowariantnej $\nabla_s$ dla koneksji Levi-Civity.
Linearyzacja równania gradientowego jest zadana wzorem
\[
 D_\gamma\xi=\nabla_s\xi+
 \nabla_\xi(\nabla_g f)\big|_{\gamma(s)}.
\]
Po wybraniu trywializacji równoległej wiązki
$\gamma^*TM$ można utożsamić ten operator z
$D_\gamma:W^{1,2}(\R,\R^d)\to L^2(\R,\R^d)$,
w którym $\nabla_s\xi$ przyjmuje postać
zwykłej pochodnej $\dd\xi/\dd s$.
W dowolnej innej trywializacji trzeba uwzględnić
dodatkową macierz współczynników koneksji.
Dla $\delta>0$ i wybranej gładkiej wagi
$\rho_\delta(s)$ zachowującej się jak
$e^{\delta|s|}$ na końcach definiujemy
$W^{1,2}_\delta$ i $L^2_\delta$ przez mnożenie przez
$\rho_\delta$. Wartość $\delta$ wybieramy mniejszą od
modułów niezerowych wartości własnych hesjanów na końcach.
\end{definicja}
\begin{uwaga}[Symetria przesunięcia]
Wektor $\dot\gamma$ jest w jądrze $D_\gamma$.
Dlatego nawet przy transwersalności nie oczekujemy
odwrotności $D_\gamma$ na całej przestrzeni.
Wprowadza się warunek cechujący, usuwający kierunek
przesunięcia czasu, oraz odpowiedni wybór dopełnień.
\end{uwaga}
\begin{lemat}[Oszacowanie błędu trajektorii sklejonej]
\label{lem:xviii-preglue}
Niech $\gamma_-:\R\to M$ i $\gamma_+:\R\to M$
będą ustalonymi, sparametryzowanymi przedstawicielami
klas w $\mathcal M(p,r)$ i $\mathcal M(r,q)$,
a ich zbieżności do $r$ niech będą wykładnicze. Dla dostatecznie dużego parametru
$R$ można skonstruować krzywą przybliżoną $\gamma_R$
powstałą przez obcięcie i sklejenie obu trajektorii,
dla której residuum
$\mathcal F(\gamma_R):=\dot\gamma_R+\nabla_g f(\gamma_R)$
spełnia $\|\mathcal F(\gamma_R)\|_{L^2_\delta}
\le C e^{-cR}$ z pewnymi stałymi $C,c>0$.
\end{lemat}
\begin{proof}
W małym otoczeniu $r$ zapisujemy fragmenty trajektorii
w ustalonej mapie wykładniczej i przesuwamy ich parametry
o $R$. Poza przedziałami, gdzie zmienia się funkcja
odcinająca, krzywa $\gamma_R$ jest dokładną trajektorią.
Na obszarze przejścia pochodna funkcji odcinającej jest
jednostajnie ograniczona, a oba sklejane fragmenty
leżą w odległości $O(e^{-\mu R})$ od $r$, dla pewnego
$\mu>0$. Gładkość $\nabla_g f$ daje tę samą kontrolę
na błąd nieliniowego członu. Po pomnożeniu przez
wagę $e^{\delta|s|}$ z $\delta<\mu$ otrzymujemy
oszacowanie $O(e^{-(\mu-\delta)R})$; stałe nie zależą
od $R$ przy ustalonych końcowych trajektoriach.
\end{proof}
\begin{lemat}[Kluczowy krok Fredholmowski --- lemat przyjmowany]
\label{lem:xviii-uniform-right-inverse}
W sytuacji Lematu~\ref{lem:xviii-preglue} i przy
transwersalności Morse'a--Smale'a wybierzmy
warunki cechujące eliminujące swobodę przesunięcia
oraz przestrzenie Banacha przekrojów
$\mathcal X_\delta=W^{1,2}_\delta$ i
$\mathcal Y_\delta=L^2_\delta$ wzdłuż $\gamma_R$.
Zakładamy zgodną z tymi wyborami konstrukcję
linearyzacji $D_R:\mathcal X_\delta\to\mathcal Y_\delta$.
Istnieje na wybranym dopełnieniu jądra operator
$Q_R:\mathcal Y_\delta\to\mathcal X_\delta$ taki, że
$D_RQ_R=\id_{\mathcal Y_\delta}$ oraz
$\|Q_R\|\le C$ dla wszystkich dostatecznie
dużych $R$, gdzie $C$ nie zależy od $R$.
\end{lemat}
\begin{uwaga}
Ten lemat nie wynika z samej małości residuum.
Aby go udowodnić, skleja się lokalne prawe odwrotne,
kontroluje błąd komutatora z funkcjami odcinającymi i
stosuje argument wykluczający koncentrację przy długiej szyjce.
Wymaga to pełnej analizy operatorów Fredholma w normach
ważonych. W tej iteracji formułujemy go precyzyjnie,
lecz \emph{nie przedstawiamy jako udowodnionego}.
\end{uwaga}
\begin{twierdzenie}[Wniosek o sklejaniu, zależny od lematu]
\label{thm:xviii-gluing-conditional}
Przy założeniach powyżej i po przyjęciu
lematu~\ref{lem:xviii-uniform-right-inverse} dla dostatecznie
dużych $R$ istnieje w pobliżu $\gamma_R$ dokładna
trajektoria $\widetilde\gamma_R$ równania gradientowego.
Jej poprawka jest rzędu $O(e^{-cR})$ w odpowiedniej normie
ważonej. Przy dodatkowym twierdzeniu o lokalnej jednoznaczności
i zgodności orientacji konstrukcja daje współrzędną końca
skompaktyfikowanej przestrzeni trajektorii.
\end{twierdzenie}
\begin{proof}
Poprawkę $\xi$ traktujemy jako mały przekrój
$\gamma_R^*TM$ i szukamy prawdziwej trajektorii
$\exp_{\gamma_R}(\xi)$, gdzie wykładnik
Riemannowski jest stosowany punktowo.
Przenosząc resztę równania równolegle wzdłuż
krótkich geodezyjnych z $\exp_{\gamma_R}(\xi)$
do $\gamma_R$, otrzymujemy odwzorowanie
Banachowskie $\mathcal F_R$ o wartościach
w stałej przestrzeni $\mathcal Y_\delta$.
Ma ono rozwinięcie
\[
\mathcal F_R(\xi)=e_R+D_R\xi+N_R(\xi),
\qquad e_R=\mathcal F_R(0),\quad
D_R=D\mathcal F_R(0),
\]
gdzie $N_R(0)=0$ i $DN_R(0)=0$.
Gładkość gradientu na zwartej rozmaitości,
jednowymiarowe oszacowanie Sobolewa
i jednostajne ograniczenia użytych funkcji
odcinających zapewniają dla małych
$\xi,\eta$ oszacowanie
$\|N_R(\xi)-N_R(\eta)\|
\le K(\|\xi\|+\|\eta\|)\|\xi-\eta\|$
ze stałą $K$ niezależną od dużego $R$.
Na wybranym dopełnieniu jądra, na które
odwzorowuje $Q_R$, równanie
$\mathcal F_R(\xi)=0$ ma postać punktu stałego
$\xi=-Q_R(e_R+N_R(\xi))$.
Istotnie, po przyłożeniu $D_R$ dostajemy
$D_R\xi=-(e_R+N_R(\xi))$
dzięki tożsamości $D_RQ_R=\id$;
odwrotnie, na wybranym dopełnieniu
$Q_RD_R=\id$, więc obie postacie są równoważne.
Lemat zapewnia $\|Q_R\|\le C$, a poprzedni lemat
$\|e_R\|\le C_0e^{-cR}$. Dla promienia
$2CC_0e^{-cR}$ prawa strona zachowuje kulę i jest
kontrakcją, gdy $R$ jest duże. Zatem istnieje poprawka
$\xi_R$ tej wielkości. Aby wniosek dał pełny opis
brzegu przestrzeni moduli, trzeba osobno wykazać, że
wszystkie bliskie trajektorie powstają taką metodą i że
orientacje są zgodne. Tych ostatnich twierdzeń w tym
dowodzie nie przyjmujemy za automatyczne.
\end{proof}

\section{Podsumowanie zakresu i zależności dowodowych}
W tym rozdziale bez zewnętrznych głębokich twierdzeń
wyprowadziliśmy definicje $\varprojlim^1$, kryterium
Mittaga--Lefflera w sytuacji przeliczalnych układów,
przykład $\Z_p/\Z$ i ciąg Milnora dla stopniowo surjektywnych
układów kompleksów. Omówiliśmy konstrukcję
wyższych przeszkód, podkreślając zależność od lokalnych
współczynników i wyboru częściowych rozszerzeń.
Dla operacji Steenroda rachunki opierają się na
przyjętym twierdzeniu Adema; dla Freudenthala na
przyjętym Blakersie--Masseyu.
W części Morse'a udowodniliśmy oszacowanie residuum i
krok kontrakcyjny warunkowo. \emph{Jednostajny prawy odwrotny,
lokalna pełność parametrów sklejania i znaki orientacji}
pozostają precyzyjnie wskazanymi tematami do dalszego dowodu.

\appendix
\renewcommand{\chaptermark}[1]{\markboth{Dodatek \thechapter}{}}
